\section{Conclusion}
\label{sec:conclusion}

Sampling cost matters. For configurable-system performance prediction, the
number of measured configurations is a useful summary but not a resource budget
unless checked against measured cost: across the systems we studied,
$n$ almost always orders measured cost, yet it can misstate the budget, hide
strategy-specific bias, and change which strategy or strategy--budget combination
is recommended.

We quantify this with the replication studies' own test. Rerunning their
Kruskal--Wallis best-sampler comparison at a fixed measured budget rather than a
fixed sample size changes the set of statistically tied-best samplers in
$38$--$55\%$ of comparisons across the five non-Linux families (and $100\%$ on
\StudyLinux), most reproducibly on the replications and Linux. The measured-cost
Pareto-efficient choice set, the stabler decision-level evidence, robustly differs
in $73$--$100\%$ of scenarios. This is a validity warning, not a systematic rejection: count-based
effort is defensible for homogeneous systems and cost-unaware samplers (we identify
several), and unsafe under heavy-tailed costs, target--cost asymmetry,
workload-sensitive distributions, dataset shifts, and diversity or
active-learning strategies that drift toward costly regions.
%: conditions common,
%measurable, and severe enough to change winners and fronts.

The stronger lesson is effectiveness: a sampler can look effective by being
accurate at a small \emph{count} budget while silently spending a large
\emph{measured} one, or look unattractive while actually costing less. Sample
count alone therefore cannot establish which strategy is cost-effective. It also
reveals an opportunity: samplers that deliberately target cheap-to-measure
configurations, and measuring cost is the first step toward designing them.

\smallskip\noindent\textbf{Open science.} Data, configurations, seeds, costs, and
reproduction scripts are in the anonymized artifact~\cite{artifact}.

\smallskip\noindent\textbf{Use of AI Assistants.} Claude Code was carefully used to assist with code
implementation and paraphrasing during paper writing. All scientific content,
experimental design, analysis, and claims are entirely the authors' own.
